\documentclass[a4paper,12pt]{article}

\usepackage[T1]{fontenc}
\usepackage[italian]{babel}
\usepackage{graphicx}
\usepackage{geometry}
\usepackage{tabularx}
\usepackage[table]{xcolor}
\usepackage[most]{tcolorbox}
\usepackage{fancyhdr}
\usepackage[hidelinks]{hyperref}

\newcommand{\groupmail}{alittlebyte.swe@gmail.com}
\newcommand{\grouplogo}{assets/logo_alittlebyte.png}
\newcommand{\groupsymbol}{assets/solo_logo.png}

\geometry{
    left=2.5cm,
    right=2.5cm,
    top=2.5cm,
    bottom=2.5cm,
    headheight=331pt,
    includeheadfoot
}

\pagestyle{fancy}
\fancyhf{}

\renewcommand{\headrulewidth}{0pt}
\fancyhead{}

\renewcommand{\footrulewidth}{0.4pt}
\fancyfoot[L]{\includegraphics[height=0.6cm]{\groupsymbol}\hspace{0.45em}aLittleByte}
    \fancyfoot[C]{\thepage}
\fancyfoot[R]{Verbale 2026-05-04}

\fancypagestyle{firstpage}{
    \fancyhf{}
    \renewcommand{\headrulewidth}{0pt}
    \renewcommand{\footrulewidth}{0.4pt}
    \fancyhead[C]{%
        \makebox[\textwidth][c]{%
            \begin{tabular}{c}
                \includegraphics[height=10cm]{\grouplogo}\\[1ex]
                \small\texttt{\groupmail}
            \end{tabular}%
        }
    }
    \fancyfoot[L]{\includegraphics[height=0.6cm]{\groupsymbol}\hspace{0.45em}aLittleByte}
        \fancyfoot[C]{\thepage}
    \fancyfoot[R]{Verbale 2026-05-04}
}

\providecommand{\glslink}[2]{\href{http://127.0.0.1:8765/glossary-html\#gls-#1}{\underline{#2}\textsuperscript{\scriptsize G}}}

\begin{document}

\thispagestyle{firstpage}

\vspace*{0.5cm}

\begin{center}
    {\LARGE \textbf{Verbale Interno - 2026-05-04}}
\end{center}

\vspace{1.5cm}

\begin{center}
    \renewcommand{\arraystretch}{1.3}
    \begin{tcolorbox}[
        width=0.92\textwidth,
        colback=gray!6,
        colframe=gray!45,
        boxrule=0.5pt,
        arc=2mm,
        boxsep=0pt,
        left=8pt, right=8pt, top=8pt, bottom=8pt
    ]
        \begin{tabularx}{\linewidth}{>{\bfseries}p{3.5cm} X}
            Gruppo      & aLittleByte (gruppo 19) \\
            Redattore   & Alex Oloni \\
            Verificatori & V. Y. Kaneda Fernandes\\
            Destinatari & Prof. Tullio Vardanega, Prof. Riccardo Cardin \\
            Data        & 04 Maggio 2026\\
        \end{tabularx}
    \end{tcolorbox}
\end{center}

\newpage

\newgeometry{left=2.5cm,right=2.5cm,top=2.5cm,bottom=2.5cm,includefoot}
\tableofcontents
\newpage
\section{Informazioni Generali}
\section*{Specifiche Documento}
\hrule
\vspace{1cm}

\begin{tabular}{>{\bfseries}p{5cm} | p{10cm}}
Luogo Riunione & \glslink{discord}{Discord} \\[6pt]

Orario (Inizio - Fine) & 15:00 -- 16.00 \\[6pt]

\glslink{redattore}{Redattore} &
  A. Oloni\\[6pt]

\glslink{amministratore}{Amministratore} & 
D. Belluz\\[6pt]

\glslink{verificatore}{Verificatore} &
V. Y. Kaneda Fernandes\\[6pt]

Partecipanti &
  A. Oloni\\ &
  L. Soligo\\&
  A. Gastaldello\\&
  A. Maule\\&
  V. Y. Kaneda Fernandes\\&

\end{tabular}


\section{Ordine del Giorno}
    \begin{itemize}
    \item Approvazione e integrazione immagini nel documento \glslink{analisi-dei-requisiti-adr}{ADR}  
    \item Gestione dei Requisiti (Opzionali e Proof of Concept)  
    \item Presentazione del foglio Excel automatizzato per il \glslink{consuntivo}{Consuntivo}  
    \item Revisione del \glslink{piano-di-progetto}{Piano di Progetto} e distribuzione ore per ruolo  
    \item \glslink{workflow}{Workflow} su \glslink{github}{GitHub} e grafici automatici
    \end{itemize}


\section{Attività svolte}
\subsection{Documentazione Analisi dei Requisiti (ADR)}
Il team ha quasi terminato l'inserimento di tutte le immagini e degli schemi necessari nel documento \glslink{analisi-dei-requisiti-adr}{ADR}. È stata identificata una sezione finale sui requisiti ancora da integrare. È emersa una discussione anche sulla mancanza di \glslink{requisiti-opzionali}{requisiti opzionali}. Si è deciso di procedere inserendo solo quanto strettamente necessario per ora, riservandosi di aggiungere ulteriori specifiche dopo un confronto con l'azienda durante la fase di PoC (Proof of Concept).

\subsection{Strumenti di Monitoraggio e Piano di Progetto}
È stato mostrato un nuovo foglio Excel automatizzato per la gestione economica: il sistema permette di aggiornare i costi e le ore in uno \glslink{sprint}{Sprint} (es. \glslink{sprint}{Sprint} 1) riflettendo automaticamente le modifiche sul saldo rimanente e sulle proiezioni degli \glslink{sprint}{Sprint} futuri. Il \glslink{piano-di-progetto}{Piano di Progetto} è stato esteso per includere anche il periodo di \glslink{pianificazione}{pianificazione} antecedente alla candidatura e le attività legate all'RTP. Si è ribadita la necessità di chiudere le \glslink{issue}{Issue} su \glslink{github}{GitHub} correttamente, poiché i diagrammi di Gantt e i report "Insights" vengono generati automaticamente per essere poi incollati nel \glslink{piano-di-progetto}{Piano di Progetto}.

\subsection{Bilancio delle Ore e dei Ruoli}
Si è discusso della ripartizione del carico di lavoro tra i vari ruoli per ottimizzare il budget:
\begin{itemize}
    \item \glslink{analista}{Analista}: Le ore preventivate risultano molto sollecitate in questa fase .
    \item \glslink{progettista}{Progettista}: Si è proposto di spostare alcune ore da questo ruolo a quello di \glslink{analista}{Analista} per compensare l'attuale carico dell'\glslink{analisi-dei-requisiti-adr}{ADR}, prevedendo che il lavoro del \glslink{progettista}{progettista} si intensificherà solo dopo l'\glslink{rtb-requirements-and-technology-baseline}{RTB}.
    \end{itemize}


\newpage
\section{To-Do List}
\begin{table}[h]
    \centering
    \begin{tabular}{|p{4.5cm}|p{5cm}|l|}
    \hline
        \rowcolor{lightgray}\textit{\textbf{Persona Incaricata}}  & \textbf{Task Assegnata} &\textbf{Link \glslink{issue}{Issue}} \\
    \hline
        \textit {Leonardo Soligo} &  Presentazione del Mock-up per l'azienda & \href{https://github.com/aLittleByte-19/Documentazione/issues/74}{\#74}\\
    \hline
        \textit{Angelica Gastaldello e Angela Maule} &  Inizio studi Backend in \glslink{laravel}{Laravel} & \href{https://github.com/aLittleByte-19/Documentazione/issues/82}{\#82}\\
    \hline
        \textit{V. Y. Kaneda Fernandes} & Continuazione della stesura del \glslink{glossario}{Glossario} & \href{https://github.com/aLittleByte-19/Documentazione/issues/53}{\#53}\\
    \hline
        \textit{Eleonora Bellet e Diego Belluz} & Inizio studi frontend in \glslink{filament}{Filament} & \href{https://github.com/aLittleByte-19/Documentazione/issues/83}{\#83}\\
    \hline 
        \textit{Alex Oloni} & Inizio stesura del \glslink{piano-di-qualifica}{Piano di Qualifica} &
        \href{https://github.com/aLittleByte-19/Documentazione/issues/92}{\#92}\\
    \hline
        \textit{Alex Oloni} & Continuazione della stesura del \glslink{piano-di-progetto}{Piano di Progetto} &
        \href{https://github.com/aLittleByte-19/Documentazione/issues/63}{\#63}\\
    \hline
        \textit{Alex Oloni} & Continuazione della stesura del File Excel per le ore & \href{https://github.com/aLittleByte-19/Documentazione/issues/70}{\#70}\\
    \hline 
        \textit{Alex Oloni} & Stesura \glslink{verbale-interno}{Verbale Interno} 04/05/2026 & \href{https://github.com/aLittleByte-19/Documentazione/issues/78}{\#78}\\
    \hline
    \end{tabular}
    \label{tab:placeholder}
\end{table}

\end{document}